\subsection{Toward a Fundamental SBC Memory Equation}

The observationally preferred SBC/F3 sector suggests that the effective infrared deformation may arise from a deeper memory dynamics.

Let \(\Psi(t)\) denote the effective observable cosmological state within a finite observational window. In the SBC framework, the observed state is not assumed to be fully local in time, but may depend on a causal memory functional,
\begin{equation}
\mathcal{M}(t)
=
\int_{-\infty}^{t}
K(t,t')\,\Psi(t')\,dt',
\end{equation}
where \(K(t,t')\) is a causal memory kernel.

Causality requires
\begin{equation}
K(t,t')=0
\qquad
{\rm for}
\qquad
t'>t .
\end{equation}

A stable memory sector must also satisfy finite normalization,
\begin{equation}
\int_{-\infty}^{t}
|K(t,t')|\,dt'
<\infty ,
\end{equation}
ensuring that remote past states do not generate divergent contributions.

The simplest kernel satisfying causality, normalization, and monotonic memory loss is
\begin{equation}
K(t,t')
=
K_0
\exp[-\lambda(t-t')]
\Theta(t-t'),
\end{equation}
where \(\Theta\) is the Heaviside causal step function.

This kernel obeys the first-order memory-decay equation
\begin{equation}
\frac{\partial K}{\partial t}
=
-\lambda K .
\end{equation}

Consequently, the effective memory amplitude satisfies
\begin{equation}
\frac{d\mathcal{M}}{dt}
=
-\lambda \mathcal{M}
+
\mathcal{S}(t),
\end{equation}
where \(\mathcal{S}(t)\) is a source term associated with the instantaneous perturbative state.

In the infrared regime, where \(\mathcal{S}(t)\) varies slowly compared with the memory-decay scale, the homogeneous solution dominates the leading screened response:
\begin{equation}
\mathcal{M}(t)
\propto
e^{-\lambda t}.
\end{equation}

Mapping the causal depth to a cosmological redshift variable gives an effective response of the form
\begin{equation}
F_{\rm SBC}(z)
\simeq
-\epsilon
\exp[-k_0(1+z)].
\end{equation}

This corresponds to the leading-order SBC/F3 phenomenological branch with
\begin{equation}
n\simeq 1.
\end{equation}

Thus, the empirical preference for \(n\simeq 1\) can be interpreted as the observational imprint of a first-order causal memory kernel.

The more general phenomenological form
\begin{equation}
F_{\rm SBC}(z)
=
-\epsilon
\exp\left[-(k_0(1+z))^n\right]
\end{equation}
may then be understood as a generalized memory-decay law, where deviations from \(n=1\) encode departures from the minimal first-order kernel.

The observational result \(n\simeq1\) therefore selects the simplest stable SBC memory regime.
